\section{Connection, obstruction and the six effects}\label{sec:two-examples}

We use the same small carrier for the two contrasting reachability
examples and the family count. Take \(X=\{0,1\}\) with its single
nonconstant parity probe and \(R=\{0,1,2,3\}\) with uniform mass.
The four marks have rate \(1/4\) and each reverses to itself.
Take \(\delta=1/8\), \(\kappa=1/2\), horizon two and two grades.
Let \(F_*\) denote the split \(\{0,1\}\mid\{2,3\}\).

The reset register has all entries \(1/4\). The gated register differs
only for one designated mark and the base pair \(g=h=0\). In that
conditional table replace its \(\{0,1\}\)-by-\(\{0,1\}\) submatrix by
\[
 \begin{pmatrix} 1/2 & 0\\ 0&1/2\end{pmatrix}.
\]
Both tables have the required stationary margins and reciprocity.
The conditional mass of each block of \(F_*\) is \(1/2\) for every
preparation in either table. The gated table is genuinely coupled to
the base pair, and loses precisely the two reciprocal proposals
between cells \(0\) and \(1\) at the designated mark and base pair.

\subsection{An exact edit through two zero responses}\label{sec:staged}

For the example of \Cref{sec:opening-example}, use the gated register
at the third mark. Hold the split at both grades, the route equality
between the first--second and first--third words, and the calibration
\(z=1/4\). Fix the fourth response at zero. Every point of
\eqref{eq:opening-path} is a lawful spectrum. Along the first segment
\(y=1/4\); along the last two \(x=0\).
The original route equality and all other tests therefore hold.

For an affine segment with varying response from \(a\) to \(b\),
choose an integer \(D>|b-a|/\varepsilon\) with \(D\ge1\) and sample
\(a+(b-a)k/D\) for \(0\le k\le D\). The difference in each marked
entry is bounded by the magnitude of the varying response change,
by the explicit two-state formula and the fixed rate and register.
The resulting rational edits are strictly within \(\varepsilon\).
The projected probability of retaining the parity bit at the
second mark changes from \(5/8\) to \(3/8\), conditional on the mark.
The prepared, marked contributions to these probabilities are
\(5/32\) and \(3/32\), respectively.
Throughout this construction the native support is fixed, since
\(Q_m>0\) and the register has the same support.

The held calibration requires \(w_z=0\) on every support containing
\(z\), and the reversal rows are identities. On
\(\{x,y,z\}\), the held word requires \(w_y=w_z\).
It follows that a single contraction cannot reach the core
\(\{z\}\), and that a surviving \(y\) keeps its positive sign.
The legal successive contractions remove \(x\) and then \(y\).
On the target support \(\{y,z\}\) the same second contraction
removes the negative \(y\). Thus both signed cores have the value
\(+1\) at \(z\) and zero elsewhere.
A rule that keeps the initial weight balances at all later supports
forces \(w_y=0\) at every step. Every reduct of each endpoint would
retain its different sign at \(y\). No common signed reduct is possible
under that frozen rule. This failure occurs even though the original
word probabilities are equal throughout the explicit chain.

\subsection{A pair that requires all six effects}\label{sec:omissions}

For a second catalogue on the same carrier, name the marks
\(a,b,c,d\) and let the initial responses be
\((1/4,1/4,-1/8,1/8)\).
Initially use the reset register and the indiscrete partition at both
grades. Hold the route \(aa\sim bb\) and the calibration
\(\alpha_b=1/4\). The optional license compares \(ca\) with \(da\),
and the optional publication checks \(\alpha_c=1/8\).
Both are defined in the fixed catalogues. All license subsets are
allowed by policy.

Construct the following chain of changes. Each row describes the
change from the state already attained by the preceding rows.
\begin{table}[ht]
\centering
\small
\begin{tabularx}{\linewidth}{@{}lXX@{}}
\toprule
Effect & Realizing change & Observable that forces the effect\\
\midrule
\(P_2\) & Gate the conditional register at mark \(c\) and base pair \(0,0\). &
Two reciprocal positive native proposals disappear.\\
\(P_1\) & Move \(\alpha_c\) from \(-1/8\) to \(1/8\) through zero. &
The measured \(c\)-response changes.\\
\(P_5\) & Disclose \(F_*\) at grade \(1\). &
The terminal decision becomes available.\\
\(P_4\) & Disclose \(F_*\) at grade \(0\). &
The earlier decision becomes available.\\
\(P_3\) & Add the license \(ca\sim da\). &
The named timed comparison is held.\\
\(P_6\) & Publish the calibration \(\alpha_c=1/8\). &
The named measurement can be cited.\\
\bottomrule
\end{tabularx}
\caption{One sequence of changes on the coupled carrier. Each adjacent pair has a lawful chain at every positive tolerance.}
\label{tab:omissions}
\end{table}
The gate is a cover, since reciprocity requires the two deleted
proposals to be removed together. The change of \(\alpha_c\) is
an affine move at fixed profile while no held numerical test
constrains that coordinate. It has arbitrary rational meshes.
Each partition advance is a cover, since nothing lies between the
indiscrete partition and a partition into two given blocks along
this refinement. Each catalogue change adds one name.
The new license passes because the \(c\) and \(d\) responses
agree by then and all the mark rates agree. The new publication
is exactly the measured \(c\)-response.

\begin{theorem}[Six measurable effects on one carrier]\label{thm:omission}
For each of the six effects there is an adjacent pair in
\Cref{tab:omissions} connected at every positive rational tolerance
whose connection requires that effect. Moreover the same initial
and final states of the entire table are connected at every positive
rational tolerance, and every connecting chain uses all six effects.
These statements also hold on the attained experimental readouts.
\end{theorem}
\begin{proof}
The preceding construction establishes connection. Every step of
another type preserves the listed observable in the third column.
For \(P_1\), this is the measured response, which all discrete steps
preserve even when a gate changes the conditional register.
For \(P_4\), no other type changes any earlier grade.
For \(P_5\), the observable is terminal availability.
The other three exclusions follow from support, the license set,
and the publication set, respectively. Thus each failure persists
under composition. The endpoints of the entire table differ in
all six listed observables, so all six types are required in every
connection. Readback follows from \Cref{prop:readback}.
\end{proof}

\subsection{A positive-tolerance obstruction}\label{sec:obstruction}

Use the initial profile of the second example and replace only
\(\alpha_a=1/4\) by \(-1/4\). Both endpoint laws pass the same held
tests. Their common calibration pins \(b\) at \(1/4\) and their
route equality forces \(\alpha_a^2=\alpha_b^2\). Thus \(w_b=0\)
and \(2w_a=2w_b\) on the occupied supports, so the respective
opposite signs of \(a\) survive. In particular the signed cores
differ. The obstruction is also quantitative. Each intermediate
law must pass the original route and calibration, hence its
\(a\)-response belongs to \(\{1/4,-1/4\}\). With \(N=8\) and
\(\delta=1/8\), a rewrite at tolerance
\[
 \frac{\delta(1/2)}{2N+1}=\frac1{272}
\]
has response change less than \(1/2\) by \eqref{eq:inverse-modulus}.
A discrete step fixes the response. Therefore no chain at this
positive tolerance can cross the two branches.
This contrasts with the earlier example in which a zero
response changes the active damping system without invalidating
the held tests.

\section{The minimum number of command families}\label{sec:basis}

An \emph{edit position} here is a writable part of the measured
record. For a numerical edit it is the local-law position.
For a discrete edit it is the support, the schedule entry at one
given grade, the license catalogue, or the publication catalogue.
The existence and uniqueness of the position follow from the cover
argument and the named readouts. A single invocation at a grade
overwrites that grade's decision entry and leaves every other
entry fixed. It still must be a legal cover between attained states.

\begin{definition}[Single-write families]
A family offers all lawful instances of one command with a fixed
mutation mode and a fixed writable field. A command for the
decision schedule has one grade argument whose domain is fixed
independently of the state and tolerance. It offers a single grade
write at each allowed argument. The four nonclock command modes
offer local-law edits, gates, route licenses and publications,
respectively. A catalogue is \emph{complete} at a tolerance if it
generates the full chain relation at that tolerance. It is a
\emph{basis} if it is complete and deleting any family destroys
completeness.
\end{definition}

\begin{lemma}[Live positions]\label{lem:live}
Call a position live at a positive tolerance if some nonidentity
elementary edit there is lawful. A catalogue is complete exactly
when it offers every live position. It is a basis exactly when
it offers all live positions and each family has a live position
offered by no other family.
\end{lemma}
\begin{proof}
Coverage is sufficient because each nonidentity elementary step
can be invoked at its position. Identity steps may be replaced
by the empty chain. For necessity, choose a live witness at an
omitted position. If it is the support, one grade, or one of the
catalogues, all other positions preserve that field, so that witness
cannot be realized. For a nonidentity local-law edit, the two
profiles agree. In any alternative chain between them, monotonicity
of support prevents a strict gate. With the local-law position
also omitted, every remaining edit fixes the law. The local
witness is again impossible. Apply the coverage test after deleting
each family to obtain the last characterization.
\end{proof}

\begin{theorem}[Five families and the terminal split]\label{thm:basis}
One uniform grade-indexed family together with the four nonclock
families offers the same elementary relation as the earlier--terminal
split, for every instrument in the class and every tolerance.
On the coupled two-grade instrument of \Cref{sec:omissions}, at
every positive rational tolerance, there are exactly two irredundant
catalogues in the declared language. These are the uniform
five-family catalogue and the split six-family catalogue.
The uniform five-family catalogue is the unique minimum-size
complete catalogue there. Five is the minimum number sufficient
for the entire class in this language.

If a successful invocation's family name alone must determine
whether it changes the terminal decision readout, the split
catalogue is the unique irredundant catalogue satisfying that
contract on the two-grade instrument.
\end{theorem}
\begin{proof}
The uniform grade family offers a decision write at every grade.
In the split catalogue the same invocation is offered by the
earlier family or by the terminal family according to its grade.
This is a bijection of invocations preserving the grade,
measured endpoints, and legality. Every finite chain is
therefore preserved.

On the coupled example all six positions are live at every
positive tolerance. The discrete examples are covers.
The response rewrite supplies a nonidentity step at every
positive tolerance by subdivision.
\Cref{lem:live} forces the four nonclock modes and coverage
of both grade positions. A family offering both grades makes
any other grade family redundant. If there is no such family,
the two singleton domains are necessary. The resulting two
catalogues are irredundant by their private live positions.
This proves the classification and the lower bound. The uniform
construction realizes that bound class-wide.

In the split catalogue every family name specifies whether the
terminal decision changes. The two live writes at the two
grades in the uniform family have different answers to that test.
The uniform catalogue therefore fails the added contract.
The classification leaves the split catalogue as the unique
irredundant one satisfying it on this instrument.
\end{proof}

There are six independently witnessed effects in
\Cref{thm:omission} and five families in the specified minimum.
Two grade-specific effects share one invocation rule while remaining
distinct measurements. The additional contract places that
distinction in the family name itself.
